\section*{v2.2.0 Structural Re-Expansion Governance Core}
\addcontentsline{toc}{section}{v2.2.0 Structural Re-Expansion Governance Core}
\begin{quote}
\textbf{Normalization note.} This block is retained as a historical governance precursor. It contributes release and re-expansion discipline, but in the normalized architecture it is not an independent proof center; proof-bearing content is carried by the grouped unfolding phases and their semantic/tensor transport rules.
\end{quote}

This release adds a governance layer to the companion monolith. Earlier stages
made closure transport, probe filtration, auditability, and ledger discipline
explicit. Version~v2.2.0 now asks an additional question: under which
rules may a compressed statement be enlarged again without damaging the closure
class, the appendix interfaces, or the downstream readout status? The answer is
cast as a governance core. Re-expansion is no longer only possible; it is to be
performed under named admissibility rules.

\subsection*{1. Governance as a proof-discipline layer}
A mature structure-first document should not merely accumulate more lemmas. It
should specify which later thickening operations are allowed and which ones are
not. Governance therefore means that every future expansion step is judged by
its compatibility with closure ancestry, probe traces, ledger surfaces, and the
status labels already attached to the claim being enlarged.

\paragraph{Governance Proposition A (admissible thickening preserves status).}
A later elaboration is admissible only if it preserves the structural status of
the compressed claim or improves it in a controlled way. An expansion that
forces a change of admissibility class without explicit declaration should be
regarded as a separate branch rather than a faithful thickening of the original
path.

\paragraph{Proof sketch.}
Compression is useful only because it abbreviates a path already licensed by the
monolith. When a later reader expands the statement again, the same path should
remain visible. If the elaboration silently changes the closure ancestry, probe
selection, or projection label, then it is not a recovery of the earlier path
but a new construction. Governance marks this distinction explicitly.

\subsection*{2. The re-expansion governance ladder}
The practical ladder of this release is
\[
\begin{aligned}
&\text{closure ancestry} \Longrightarrow \text{probe/interface trace} \Longrightarrow \text{ledger or appendix encoding}\\
&\Longrightarrow \text{compressed readout} \Longrightarrow \text{governed re-expansion}\\
&\Longrightarrow \text{status check} \Longrightarrow \text{return to the ledger}.
\end{aligned}
\]
The final return arrow is part of the new emphasis. Re-expansion is not treated
as a one-way release of detail. It must return to the ledger surface in a form
that remains compatible with the proof memory of the monolith.

\paragraph{Governance Proposition B (every governed expansion returns to a ledger surface).}
A mathematically thickened block is fully integrated only when its added detail
can be pushed back into a stable appendix, tool, lemma, or table interface.
Otherwise the expansion remains locally useful but globally insufficiently
synchronized.

\subsection*{3. Probe governance and admissible branch formation}
The probe channel now gains a sharper branching role. Not every elaboration of a
probe-selected mode has to remain in the same local branch, but any branch split
must be declared at the point where the probe trace ceases to determine a unique
downstream continuation. This keeps the search-space logic compatible with the
proof memory logic.

\paragraph{Governance Proposition C (branching must occur at visible trace loss).}
If a later derivation no longer preserves a unique recognizable probe trace,
then a branch declaration is required. Continuing as if the same path remained
in force would obscure which downstream statements are still inherited and which
are newly introduced.

\subsection*{4. Appendix governance and reversible table discipline}
The appendix receives a correspondingly stronger role. Rows, tables, and named
tools are now governed interfaces. They should not only compress results and
carry reopening metadata; they should also state which kinds of later expansion
may legitimately pass through them. In this sense the appendix becomes a modest
control layer for future mathematical thickening.

\paragraph{Governance Proposition D (table rows define admissible recovery windows).}
A mature table row should delimit the window inside which a later elaboration is
still regarded as a recovery of the same path class. Once an expansion leaves
that window, the result should either be relabeled or entered as a new row,
lemma, or branch.

\subsection*{5. Consequence for the next proof cycle}
With v2.2.0 the companion monolith now carries not only closure,
projection, ledger, and audit discipline, but also a clearer rule for how later
mathematical detail may be inserted without dissolving the structure-first
architecture. The next proof cycle should therefore thicken selected paths under
explicit governance rather than by unconstrained local expansion. This is the
cleanest route toward the larger companion target with many more lemmas and
technical proof tools.

\paragraph{Working rule for subsequent versions.}
Whenever a new technical derivation is inserted, it should be possible to say
(i) which closure and probe ancestry it inherits, (ii) which appendix or ledger
surface compresses it, (iii) which re-expansion window legitimizes its added
detail, and (iv) how the resulting thickening returns to a synchronized proof
interface. Material meeting all four demands strengthens the monolith in the
intended disciplined way.


\clearpage
